\documentclass[11pt,a4paper]{article}

\usepackage[spanish]{babel}
\usepackage[utf8]{inputenc}
\usepackage[T1]{fontenc}
\usepackage{amsmath}
\usepackage{amssymb}
\usepackage{xcolor}
\usepackage{enumitem}
\usepackage{geometry}
\usepackage{listings}
\usepackage{hyperref}

\geometry{margin=2.5cm}

\definecolor{codegray}{rgb}{0.95,0.95,0.95}

\lstset{
backgroundcolor=\color{codegray},
basicstyle=\ttfamily\small,
frame=single,
breaklines=true,
columns=fullflexible
}

\title{\textbf{Proyecto Funcional}\\ Procesamiento Paralelo de Imágenes con Scheme y Erlang}
\author{Eddy Ramírez Jiménez}
\date{\today}

\begin{document}

\maketitle

\section{Descripción general}

El objetivo del proyecto es desarrollar un sistema para el procesamiento
paralelo de imágenes que integre los lenguajes de programación funcional
\textbf{Scheme} y \textbf{Erlang}.

El sistema deberá utilizar Erlang como lenguaje encargado de la
concurrencia y coordinación del procesamiento, mientras que Scheme será
responsable de realizar las transformaciones sobre las imágenes.

La arquitectura general del sistema será:

\begin{center}
\begin{verbatim}
+----------------+
|     Erlang     |
|   Coordinador  |
+-------+--------+
        |
Lee y divide la imagen
        |
+-------+---------+-----------------+
|                 |                 |
v                 v                 v
+----------+      +----------+      +----------+
| Proceso  |      | Proceso  |      | Proceso  |
| Erlang 1 |      | Erlang 2 |      | Erlang 3 |
+----+-----+      +----+-----+      +----+-----+
     |                 |                 |
     v                 v                 v
  Scheme            Scheme            Scheme
     |                 |                 |
+----+-----------------+-----------------+
|
v
Imagen reconstruida
\end{verbatim}
\end{center}

Cada proceso Erlang deberá encargarse de una región de la imagen y
deberá ejecutar una instancia independiente de Scheme para procesarla.

\section{Objetivos}

Al finalizar el proyecto, el estudiante deberá demostrar dominio de los
siguientes conceptos:

\begin{itemize}
\item Programación funcional mediante Scheme.
\item Manipulación recursiva de estructuras de datos.
\item Representación y procesamiento de imágenes.
\item Creación y comunicación entre procesos Erlang.
\item Paralelización de tareas.
\item Comunicación entre programas escritos en diferentes lenguajes.
\item Diseño de protocolos de comunicación.
\item Medición y análisis del rendimiento de programas paralelos.
\end{itemize}

\section{Especificación del sistema}

El sistema recibirá como entrada una imagen y una cantidad de procesos
concurrentes. La imagen deberá ser procesada mediante un filtro
implementado en Scheme.

Como mínimo, el sistema deberá implementar un \textbf{difuminado
gaussiano (Gaussian Blur)}.

La ejecución deberá seguir, conceptualmente, las siguientes etapas:

\begin{enumerate}
\item Erlang recibe el nombre de la imagen de entrada.
\item Erlang lee la imagen.
\item Erlang divide la imagen en regiones.
\item Erlang crea los procesos encargados del procesamiento.
\item Cada proceso recibe una región de la imagen.
\item Cada proceso ejecuta una instancia independiente de Scheme.
\item Scheme procesa la región utilizando el filtro correspondiente.
\item Scheme devuelve la región procesada.
\item Erlang recibe los resultados.
\item Erlang reconstruye la imagen original a partir de las regiones
procesadas.
\item Erlang genera la imagen de salida.
\end{enumerate}

\section{Formato de imagen}

Para evitar que el proyecto dependa de bibliotecas externas de
procesamiento de imágenes, se deberá utilizar el formato
\textbf{PPM}.

Se recomienda utilizar el formato P3, cuya estructura básica es:

\begin{lstlisting}
P3
800 600
255
255 0 0
255 0 0
...
\end{lstlisting}

El programa deberá ser capaz de leer y escribir correctamente imágenes
PPM.

Cada píxel deberá poder representarse mediante sus componentes:

$$
(R,G,B)
$$

donde cada componente representa un valor entre $0$ y $255$.

\section{Parte I: procesamiento en Scheme}

Scheme será responsable exclusivamente del procesamiento de una región
de la imagen.

El programa en Scheme deberá implementar, como mínimo, una función
equivalente a:

\begin{lstlisting}
(gaussian image-region kernel)
\end{lstlisting}

La función deberá recibir una región de la imagen y un kernel de
convolución y deberá producir una nueva región con el filtro aplicado.

Como mínimo, deberá soportarse un kernel gaussiano de tamaño $3\times3$.

Por ejemplo:

$$
G=
\frac{1}{16}
\begin{pmatrix}
1&2&1\\
2&4&2\\
1&2&1
\end{pmatrix}.
$$

Para cada píxel, el filtro deberá calcular una combinación ponderada
de los píxeles vecinos.

El procesamiento deberá realizarse utilizando las características
propias de Scheme. Se espera el uso apropiado de funciones, recursión,
listas y abstracciones funcionales.

\subsection{Restricciones sobre Scheme}

Scheme no deberá encargarse de:

\begin{itemize}
\item Crear los procesos paralelos.
\item Distribuir el trabajo.
\item Coordinar los procesos.
\item Reconstruir la imagen completa.
\end{itemize}

Estas responsabilidades corresponden a Erlang.

\section{Parte II: procesamiento en Erlang}

Erlang será el responsable de coordinar toda la ejecución.

Deberá implementar un proceso coordinador que:

\begin{enumerate}
\item Lea la imagen de entrada.
\item Determine sus dimensiones.
\item Divida la imagen en regiones.
\item Cree los procesos necesarios.
\item Distribuya las regiones entre los procesos.
\item Ejecute Scheme para cada región.
\item Reciba los resultados.
\item Reconstruya la imagen.
\item Escriba la imagen resultante.
\end{enumerate}

El número de procesos deberá poder especificarse como parámetro de
ejecución.

Por ejemplo:

\begin{lstlisting}
./image_processor imagen.ppm salida.ppm 4
\end{lstlisting}

indicaría que la imagen deberá procesarse utilizando cuatro procesos
Erlang.

\section{Comunicación entre Erlang y Scheme}

Uno de los objetivos fundamentales del proyecto es integrar dos
lenguajes diferentes.

Cada proceso Erlang deberá ejecutar una instancia independiente del
programa Scheme.

El mecanismo concreto de comunicación entre ambos lenguajes deberá ser
diseñado e implementado por los estudiantes.

El protocolo deberá permitir, como mínimo:

\begin{itemize}
\item Enviar una región de imagen a Scheme.
\item Especificar los parámetros del filtro.
\item Indicar el tamaño del kernel.
\item Recibir la región procesada.
\item Detectar errores durante la ejecución de Scheme.
\end{itemize}

El informe deberá explicar claramente el protocolo utilizado y
justificar las decisiones de diseño.

\section{Procesamiento de regiones}

La imagen deberá dividirse en regiones para permitir el procesamiento
concurrente.

Por ejemplo, una imagen podría dividirse de la siguiente manera:

\begin{center}
\begin{verbatim}
+-------------------+-------------------+
|                   |                   |
|        R1         |        R2         |
|                   |                   |
+-------------------+-------------------+
|                   |                   |
|        R3         |        R4         |
|                   |                   |
+-------------------+-------------------+
\end{verbatim}
\end{center}

El método exacto de división queda a criterio del estudiante.

Sin embargo, los filtros de convolución presentan un problema
importante: el valor de un píxel puede depender de píxeles ubicados en
regiones vecinas.

Por ejemplo, para un kernel $3\times3$, un píxel cercano al límite entre
dos regiones puede necesitar información de la región contigua.

Por esta razón, el sistema deberá implementar un mecanismo de
\textbf{halo} o región de contexto.

Cada proceso podrá recibir píxeles adicionales alrededor de su región,
pero deberá devolver únicamente los píxeles correspondientes a la
región que le fue asignada.

\section{Procesamiento de los bordes}

El sistema deberá especificar qué sucede cuando el kernel requiere
píxeles que se encuentran fuera de la imagen.

Se permitirá utilizar cualquiera de las siguientes estrategias:

\begin{itemize}
\item Extender el borde.
\item Utilizar píxeles con valor cero.
\item Repetir el píxel más cercano.
\item Reflejar la imagen.
\end{itemize}

La estrategia seleccionada deberá estar documentada y ser utilizada
consistentemente.

\section{Concurrencia}

La solución deberá utilizar el modelo de concurrencia de Erlang.

Cada región deberá ser procesada por un proceso Erlang independiente.

La comunicación entre procesos deberá realizarse mediante el mecanismo
de paso de mensajes de Erlang.

No se permitirá implementar la paralelización mediante una simple
secuencia de llamadas.

La solución deberá demostrar que las regiones pueden procesarse de
forma concurrente.

\section{Tolerancia a fallos}

El sistema deberá ser capaz de detectar el fallo de un proceso
encargado de procesar una región.

Como mínimo, deberá detectarse cuando una instancia de Scheme termine
incorrectamente.

El sistema deberá informar el error y evitar que una falla silenciosa
produzca una imagen incorrecta.

Como extensión, se recomienda implementar la posibilidad de volver a
ejecutar automáticamente una región cuyo procesamiento haya fallado.

\section{Modos de ejecución}

El sistema deberá permitir comparar al menos dos configuraciones:

\begin{enumerate}
\item Ejecución utilizando un único proceso.
\item Ejecución utilizando múltiples procesos.
\end{enumerate}

Se recomienda realizar experimentos con:

$$
1,\;2,\;4,\;8
$$

procesos, cuando el hardware disponible lo permita.

Para cada configuración deberán registrarse los tiempos de ejecución.

\section{Evaluación del rendimiento}

El informe deberá analizar el comportamiento del sistema utilizando
como mínimo las siguientes métricas.

\subsection{Tiempo de ejecución}

Para cada cantidad de procesos $p$, se deberá medir:

$$
T_p
$$

correspondiente al tiempo necesario para procesar la imagen completa.

\subsection{Speedup}

Se deberá calcular:

$$
S_p=\frac{T_1}{T_p}.
$$

\subsection{Eficiencia}

También deberá calcularse:

$$
E_p=\frac{S_p}{p}.
$$

Los resultados deberán presentarse mediante tablas y gráficos.

El análisis deberá explicar las razones por las cuales el speedup
obtenido puede ser menor que el speedup ideal.

Entre los aspectos que deberán considerarse se encuentran:

\begin{itemize}
\item Creación de procesos.
\item Comunicación entre procesos.
\item Comunicación entre Erlang y Scheme.
\item Lectura y escritura de archivos.
\item División y reconstrucción de la imagen.
\item Copia y serialización de datos.
\end{itemize}

\section{Extensiones}

Además del filtro gaussiano obligatorio, se deberá implementar al menos
\textbf{una transformación adicional}.

Algunas posibilidades son:

\begin{itemize}
\item Escala de grises.
\item Inversión de colores.
\item Ajuste de brillo.
\item Threshold.
\item Sharpen.
\item Detección de bordes mediante Sobel.
\end{itemize}

La transformación adicional deberá ser implementada en Scheme.

Como extensión adicional, se podrá permitir que el usuario especifique
una secuencia de transformaciones.

Por ejemplo:

\begin{lstlisting}
(pipeline
(grayscale)
(gaussian 3)
(threshold 128))
\end{lstlisting}

\section{Interfaz mínima}

El sistema deberá proporcionar una interfaz que permita ejecutar al
menos:

\begin{lstlisting}
./image_processor entrada.ppm salida.ppm 1
./image_processor entrada.ppm salida.ppm 2
./image_processor entrada.ppm salida.ppm 4
./image_processor entrada.ppm salida.ppm 8
\end{lstlisting}

El último parámetro representa la cantidad de procesos Erlang que
deberán utilizarse.

La interfaz podrá modificarse siempre que conserve la funcionalidad
solicitada.

\section{Entregables}

El proyecto deberá incluir:

\begin{enumerate}

\item \textbf{Imágenes de prueba}

Un conjunto de imágenes utilizadas para verificar el funcionamiento
del sistema.

\item \textbf{Manual de ejecución}

Instrucciones suficientes para compilar, configurar y ejecutar el
proyecto.

\item \textbf{Informe técnico}

El informe deberá explicar:

\begin{itemize}
    \item Arquitectura del sistema.
    \item División de la imagen.
    \item Modelo de concurrencia utilizado.
    \item Comunicación entre Erlang y Scheme.
    \item Protocolo de comunicación.
    \item Implementación del filtro.
    \item Manejo de los bordes.
    \item Manejo de errores.
    \item Resultados experimentales.
    \item Análisis del speedup y eficiencia.
    \item Limitaciones del sistema.
\end{itemize}

\item \textbf{Demostración}

El proyecto deberá ser presentado y ejecutado durante la
demostración correspondiente.

\end{enumerate}

\section{Restricciones}

\begin{itemize}

\item El procesamiento del filtro obligatorio deberá realizarse en
Scheme.

\item La creación y coordinación de los procesos deberá realizarse
en Erlang.

\item Cada región deberá ser procesada mediante una instancia
independiente de Scheme ejecutada desde Erlang.

\item No se permite realizar todo el procesamiento de la imagen en
un único proceso Scheme y utilizar Erlang únicamente como
lanzador.

\item No se permite utilizar una biblioteca externa que implemente
directamente el filtro gaussiano.

\item La solución deberá poder ejecutarse utilizando las versiones
de Erlang y Scheme indicadas por el profesor.

\end{itemize}

\section{Criterios de evaluación}

La calificación se distribuirá de la siguiente manera:

\begin{center}
\begin{tabular}{lr}
\hline
\textbf{Componente} & \textbf{Porcentaje}\\
\hline
Implementación en Scheme & 20\%\\
Implementación en Erlang & 25\%\\
Integración Erlang--Scheme & 20\%\\
Concurrencia y distribución & 15\%\\
Procesamiento correcto de imágenes & 10\%\\
Informe y análisis experimental & 10\%\\
\hline
\textbf{Total} & \textbf{100\%}\\
\hline
\end{tabular}
\end{center}

\section{Defensa del proyecto}

Durante la defensa, los estudiantes deberán ser capaces de explicar
cualquier parte del código desarrollado.

El evaluador podrá solicitar modificaciones pequeñas al programa durante
la defensa, por ejemplo:

\begin{itemize}
\item Cambiar el tamaño del kernel.
\item Modificar la cantidad de procesos.
\item Cambiar la estrategia para los bordes.
\item Incorporar un nuevo filtro.
\item Forzar el fallo de un proceso.
\item Modificar la forma de dividir la imagen.
\end{itemize}

La defensa deberá demostrar que los integrantes comprenden tanto la
implementación en Scheme como la implementación en Erlang.

La fecha de entrega será el día 18 de setiembre.

\end{document}